\documentclass[12pt]{article}
\usepackage[T1]{fontenc}
\usepackage[utf8]{inputenc}
\usepackage{lmodern}
\usepackage{textcomp}
\usepackage{enumitem}
\usepackage{geometry}
\geometry{margin=1in}
\begin{document}
\begin{center}
Answer Key - Philosophy - Undergraduate Level Assessment 9 \
\small (Answers are ordered exactly as in the question paper.)
\end{center}

\section*{Multiple Choice}
\begin{enumerate}[label=\arabic*.]
\item \textbf{Answer:} A
\\ \textit{Options:} A) cost.; B) weight.; C) species.; D) purity.; E) morality.
\\ \textit{Note:} Correct option: A - cost.
\item \textbf{Answer:} B
\\ \textit{Options:} A) performed outside of marriage; B) non-consensual.; C) practiced among adolescents; D) not leading to emotional fulfillment; E) not accepted by societal norms
\\ \textit{Note:} Correct option: B - non-consensual.
\item \textbf{Answer:} E
\\ \textit{Options:} A) Roman Catholic; B) Lutherans; C) Methodists; D) Greek Orthodox; E) Quakers
\\ \textit{Note:} Correct option: E - Quakers
\item \textbf{Answer:} A
\\ \textit{Options:} A) Papadaki actually agrees with Kant's view; B) Papadaki suggests that Kant's view of sexual objectification is too subjective.; C) Papadaki argues that sexual objectification is an outdated concept.; D) sexual objectification actually enhances humanity.; E) Papadaki believes that sexual objectification is a necessary part of human relationships.
\\ \textit{Note:} Correct option: A - Papadaki actually agrees with Kant's view
\item \textbf{Answer:} E
\\ \textit{Options:} A) all be avoidable.; B) only felt by multiple people at different times.; C) all be experienced in the same way.; D) all be of the same intensity.; E) all felt by a single person.
\\ \textit{Note:} Correct option: E - all felt by a single person.
\end{enumerate}

\section*{True / False}
\begin{enumerate}[label=\arabic*.]
\setcounter{enumi}{5}
\item \textbf{Answer:} True
\\ \textit{Options:} A) True; B) False
\\ \textit{Note:} LLM-generated true statement based on MCQ.
\item \textbf{Answer:} True
\\ \textit{Options:} A) True; B) False
\\ \textit{Note:} LLM-generated true statement based on MCQ.
\item \textbf{Answer:} False
\\ \textit{Options:} A) True; B) False
\\ \textit{Note:} LLM-generated false statement. Correct answer: 'can be compared, but only roughly.'.
\item \textbf{Answer:} False
\\ \textit{Options:} A) True; B) False
\\ \textit{Note:} LLM-generated false statement. Correct answer: 'is not always hate speech, because it does not involve violence'.
\item \textbf{Answer:} False
\\ \textit{Options:} A) True; B) False
\\ \textit{Note:} LLM-generated false statement. Correct answer: 'an expression of the ethics of care.'.
\end{enumerate}

\section*{Long Form Response}
\begin{enumerate}[label=\arabic*.]
\setcounter{enumi}{10}
\item \textbf{Answer:} Red herring argument. Explain why this is the correct answer and provide supporting evidence. Reference key facts or concepts that support this choice.
\\ \textit{Note:} Long-form derived from MCQ; award credit for stating the correct idea and giving supporting reasoning.
\item \textbf{Answer:} The minor premise must either accept an alternative or reject an alternative. Explain why this is the correct answer and provide supporting evidence. Reference key facts or concepts that support this choice.
\\ \textit{Note:} Long-form derived from MCQ; award credit for stating the correct idea and giving supporting reasoning.
\end{enumerate}

\vfill
\noindent\textit{End of Answer Key}
\end{document}